\documentclass[acmsmall,screen]{acmart}
%%
%% \BibTeX command to typeset BibTeX logo in the docs
\AtBeginDocument{%
  \providecommand\BibTeX{{%
    Bib\TeX}}}

\setcopyright{acmlicensed}
\copyrightyear{2026}
\acmYear{2026}
\acmDOI{XXXXXXX.XXXXXXX}

\acmJournal{JRC}
\acmVolume{0}
\acmNumber{0}
\acmArticle{0}
\acmMonth{0}

\acmSubmissionID{XXX-XXX-XXX}

\begin{document}

\title{Public interest data infrastructuring}

\author{Brian C. Keegan}
\email{brian.keegan@colorado.edu}
\orcid{0000-0002-7793-398X}
\affiliation{%
  \institution{University of Colorado Boulder}
  \city{Boulder}
  \state{Colorado}
  \country{USA}
}

\renewcommand{\shortauthors}{Keegan}

\begin{abstract}
Data-driven systems increasingly govern critical aspects of public life, yet access to the data and infrastructures needed to study and oversee these systems is rapidly shrinking. This article defines public interest data infrastructuring as an emerging field that adapts the public missions of journalism, law, accounting, planning, and engineering to contemporary data-intensive environments. It argues that public interest data infrastructuring arises in response to three structural dynamics: enclosure, where once-public data is privatized or locked down; exemption, where platforms and AI developers operate outside meaningful regulatory scrutiny; and erosion, where fragmentation and infrastructural decay undermine long-term observation and reproducible research. Against these forces, the article advances a constructive agenda organized around openness, oversight, and ownership as countervailing pillars of accountability and stewardship. Public interest data infrastructuring is distinguished from adjacent efforts such as data for good, civic technology, critical data studies, and data justice by its emphasis on building durable infrastructures and regulatory architectures for democratic oversight. Three case studies on mining municipal archives, infrastructuring worker observatories, and designing Anthropocene data trusts illustrate how public interest data infrastructuring reimagines the relationships between data, power, and the public good in practice. The article concludes by outlining how public interest data infrastructuring can provide the foundations for future institutions grounded in contemporary regulatory precedents.
\end{abstract}

\begin{CCSXML}
<ccs2012>
   <concept>
       <concept_id>10003456.10003457.10003567.10010990</concept_id>
       <concept_desc>Social and professional topics~Socio-technical systems</concept_desc>
       <concept_significance>500</concept_significance>
       </concept>
   <concept>
       <concept_id>10003456.10003462.10003480.10003484</concept_id>
       <concept_desc>Social and professional topics~Technology and censorship</concept_desc>
       <concept_significance>300</concept_significance>
       </concept>
   <concept>
       <concept_id>10003456.10003462.10003544.10003589</concept_id>
       <concept_desc>Social and professional topics~Governmental regulations</concept_desc>
       <concept_significance>300</concept_significance>
       </concept>
   <concept>
       <concept_id>10010405.10010455</concept_id>
       <concept_desc>Applied computing~Law, social and behavioral sciences</concept_desc>
       <concept_significance>500</concept_significance>
       </concept>
   <concept>
       <concept_id>10010405.10010476.10010936</concept_id>
       <concept_desc>Applied computing~Computing in government</concept_desc>
       <concept_significance>300</concept_significance>
       </concept>
    <concept>
        <concept_id>10003120.10003130.10003131</concept_id>
        <concept_desc>Human-centered computing~Collaborative and social computing~Collaborative and social computing theory, concepts and paradigms</concept_desc>
        <concept_significance>300</concept_significance>
        </concept>
    <concept>
        <concept_id>10010405.10010476.10003392</concept_id>
        <concept_desc>Applied computing~Digital libraries and archives</concept_desc>
        <concept_significance>300</concept_significance>
        </concept>
 </ccs2012>
\end{CCSXML}

\ccsdesc[500]{Social and professional topics~Socio-technical systems}
\ccsdesc[300]{Social and professional topics~Technology and censorship}
\ccsdesc[300]{Social and professional topics~Governmental regulations}
\ccsdesc[500]{Applied computing~Law, social and behavioral sciences}
\ccsdesc[300]{Applied computing~Computing in government}
\ccsdesc[300]{Applied computing~Digital libraries and archives}
\ccsdesc[300]{Human-centered computing~Collaborative and social computing~Collaborative and social computing theory, concepts and paradigms}

\keywords{data governance; information commons; accountability; transparency; governance; regulation; infrastructuring; municipal data; gig workers; climate justice}

\received{20 February 2007}
\received[revised]{12 March 2009}
\received[accepted]{5 June 2009}

\maketitle

\section*{Introduction}

For more than a century, democratic societies have depended on data to render public life visible and governable. Official statistics, regulatory filings, and administrative records have underpinned projects as diverse as labor protection, civil rights enforcement, environmental regulation, and economic planning. In each domain, new forms of measurement were accompanied by struggles over access, interpretation, and control: who could see the numbers, who could challenge them, and who was accountable for their consequences. The rise of large-scale digital platforms in the early twenty-first century appeared, for a time, to extend this evidentiary infrastructure into the everyday fabric of social interaction. 

For roughly a decade between 2008 and 2018, companies such as Twitter, Facebook, Reddit, and Google invested in APIs, data dumps, and other interfaces that created a ``golden age'' for studying social behavior at unprecedented scale and fidelity. These infrastructures functioned like scientific instruments: observatories producing high-resolution behavioral data about social structure, linguistic dynamics, cultural practices, and algorithmic interventions~\cite{Bruns_APIcalypsesocialmedia_2019,lazer_ComputationalSocialScience_2020,proferes_StudyingRedditSystematic_2021,Hampton_StudyingDigitalDirections_2017,GolderDigitalFootprintsOpportunities2014,Lazer_DataexMachina_2017}. Researchers in data science, computational social science, human–computer interaction, and related fields built methods, careers, and new subdisciplines on these platforms' data.

Platforms initially welcomed this work because it conferred legitimacy: analyses that evaluated new features, documented rapid growth, and demonstrated ``real-world impact'' made platforms more attractive to users, employees, investors, and regulators. But as these firms matured into publicly traded corporations, their incentives shifted. By the mid-2010s, researchers were using the same data infrastructures to uncover bots, harassment, polarization, disinformation, and other phenomena implicating platforms' governance failures and systemic biases~\cite{marwick_and_lewis,Walker_disinformationlandscapelockdown_2019}. In response, companies adopted technical and legal strategies that made it increasingly difficult to study their systems. Access policies tightened, terms of service hardened, and claims about protecting user privacy became convenient justifications for shutting down data access to independent investigators~\cite{confessore_cambridge_2018,bobrowsky_FacebookDisablesAccess_2021,stokel-walker_under_2024}. Freelon aptly described the resulting condition as a ``post-API age,'' in which researchers' access to digital trace data can be restricted at any time, for any reason, and without recourse~\cite{Freelon_ComputationalResearchPostAPI_2018}.

The emergence of large language models in the early 2020s accelerated this transformation. User-generated content accumulated over decades suddenly became highly valuable as training data. Platforms moved swiftly to prevent scraping, restrict APIs, and paywall archives, treating their data as proprietary assets rather than quasi-public resources~\cite{isaac_RedditWantsGet_2023,roose_data_2024}. At the same time, user migration to alternative platforms and federated networks has made even basic measurement of online social behavior more difficult~\cite{nicholson_truthiness_2024}. Few organizations outside the largest technology companies have the technical and financial capacity to sustain open data infrastructures at scale. In retrospect, the ``API age'' was not only unusually open but also riddled with inefficiencies, inequalities, and ethical abuses that platforms have seized upon to justify restricting access~\cite{boydCriticalQuestionsBig2012,Puschmann_endwildwest_2019}. The impacts are akin to new management putting the Large Hadron Collider's data behind a paywall or the James Webb Space Telescope was being abandoned because it was generating inconvenient data. 

There is no straightforward path back to the peculiar data ecosystem of the late 2000s and early 2010s. Unlike particle physicists or astronomers, computational social scientists and data scientists working in the public interest generally do not control the instruments that generate the data on which their research depends. The challenge extends beyond academic research and into the agendas of other public interest watchdogs. As data-driven systems mediate communication, work, health, mobility, and political expression, restrictions on access and scrutiny jeopardize democratic oversight more broadly. While social media platforms might prefer to keep our regulatory imaginations small, there are many historical precedents for compelling private organizations to produce data in the public interest.

This situation is not a temporary policy problem but the outcome of three deeper, mutually reinforcing forces: enclosure, exemption, and erosion. \textit{Enclosure} refers to the privatization or withdrawal of data that once had public or research value, whether through API shutdowns, paywalls, or the quiet retirement of government datasets. \textit{Exemption} denotes the ability of powerful platforms and AI developers to operate outside meaningful regulatory scrutiny, invoking privacy, intellectual property, or competitive sensitivity to avoid disclosure and accountability. \textit{Erosion} captures the decay of institutions and infrastructures like APIs and data dumps that once enabled long-term observation and reproducible analysis of digital life. Together, these forces are transforming the informational foundations of public knowledge and weakening society's capacity to govern data-intensive systems. In the language of infrastructure studies, these are also \textit{infrastructuring} traditions: they treat accountability as something that must be continuously produced through records, standards, routines, and durable institutions, not simply asserted as a value. They therefore supply tested patterns for how openness, oversight, and collective stewardship are enacted and maintained under real-world constraints of scarcity, adversarial incentives, and political pressure.

This article proposes ``public interest data infrastructuring'' as an interdisciplinary project grounded in explicit normative values that takes these structural pressures seriously and seeks to build countervailing commitments of openness, oversight, and ownership. It follows arguments in other domains that locate professions within longer lineages of public interest traditions like journalism, law, accounting, engineering, planning that have developed mechanisms to ensure that specialized expertise serves the common good rather than only private interests~\cite{washingtonDefiningPublicInterest2024}. Drawing on these lineages and on adjacent conversations in civic technology, data-for-good, critical data studies, and data justice, public interest data infrastructuring aims to reconstruct the conditions under which data-driven systems can credibly claim to operate in the public interest.

The remainder of the article proceeds in four steps. First, it offers a conceptual diagnosis of enclosure, exemption, and erosion as structural forces undermining public oversight of data-driven systems. Second, it defines public interest data infrastructuring and elaborates openness, oversight, and ownership as counter-pressures. Third, it triangulates public interest data infrastructuring with related public interest lineages from journalism, law, accounting, planning, and engineering as well as public interest conversations like data for good, civic technology, and critical data studies. Fourth, common critiques are summarized and addressed. Finally, three speculative case studies illustrate how public interest data infrastructuring might be enacted in practice: mining municipal archives, infrastructuring worker observatories, and Anthropocene data trusts. The concluding discussion synthesizes anticipated critiques, draws out implications for regulation, professional norms, and grassroots practice, and sketches a forward-looking agenda for building data futures oriented toward the public, workers, and future generations.

\section*{Three pressures}

\paragraph*{Enclosure}
The first defining force is enclosure. Beginning in medieval and early modern England, the enclosure movement transformed collectively held lands into private property by fencing off fields, revoking customary rights, and displacing communities who relied on shared resources~\cite{boyle2017second}. Enclosure was not merely an economic shift; it was a political act that redistributed power, restricted access, and redefined who could benefit from what had once been a commons~\cite{hess2003ideas}. A similar dynamic now characterizes the digital ecosystem~\cite{boyle2017second,hess2003ideas,wu_PlatformEnclosureHuman_2021}. Information enclosure refers to the privatization and withdrawal of access to data that once supported independent judgments about how digital systems operate. In rapid succession, Facebook restricted access to its public APIs after the Cambridge Analytica scandal~\cite{confessore_cambridge_2018,van_der_vlist_api_2022}, Twitter revoked research access following its 2022 acquisition~\cite{han_what_2024,stokel-walker_under_2024}, and Reddit shut down the Pushshift infrastructure while imposing steep fees for API access~\cite{isaac_RedditWantsGet_2023,lift_ticket83_RedditDataAPI_2023}. These actions mirror historical enclosure: platforms have fenced off user-generated content that previously functioned as a public or quasi-public resource for researchers, journalists, policymakers, and civil society~\cite{Freelon_ComputationalResearchPostAPI_2018,Bruns_APIcalypsesocialmedia_2019,wu_PlatformEnclosureHuman_2021}. Governments have likewise reduced or eliminated public datasets on policing, environmental quality, public health, and education.

\paragraph*{Exemption}
The second force is exemption, which describes how companies, their platforms, and their algorithms operate outside meaningful regulatory scrutiny~\cite{suzor2018digital}. As their influence over communication, commerce, and employment has grown, these firms have invoked legal and technical shields that insulate them from oversight. Companies routinely claim privacy, intellectual property, national security, or competitive sensitivity to justify withholding data and treating opacity as a competitive advantage~\cite{suzor2018digital,obar2020biggest}. They selectively invoke decades-old laws such as the Computer Fraud and Abuse Act to prevent independent data collection~\cite{ACLU_Sandvig_2019} and Section 230 to disclaim responsibility for harms facilitated by their systems. Platforms increasingly engage in ``privacy-washing'' by claiming to protect users while in practice restricting independent research and accountability~\cite{obar2020biggest,van_der_vlist_api_2022}. Even in the face of well-documented harms like bias, harassment, misinformation, discrimination, and environmental externalities, companies face minimal obligation to disclose relevant data or to mitigate harms. Exemption ensures that socio-technical systems with profound social, economic, and political impact operate with limited accountability~\cite{suzor2018digital}. It thus creates a structural mismatch: the systems with the greatest public consequences are subject to the weakest oversight.

\paragraph*{Erosion}
Erosion, the third defining force, refers to the ongoing breakdown of the institutions and infrastructures that once enabled sustained, independent observation of platform-mediated behavior~\cite{Plantin_Infrastructurestudiesmeet_2016}. As researchers lose access to centralized APIs and archival datasets, studying long-term trends in digital behavior becomes increasingly difficult~\cite{Freelon_ComputationalResearchPostAPI_2018}. The rise of decentralized and federated platforms further disperses data across fragmented, incommensurable, and ephemeral infrastructures. Basic indicators that once grounded empirical research have become harder to measure consistently~\cite{Bruns_APIcalypsesocialmedia_2019,nicholson_truthiness_2024}. Erosion undermines the reproducibility, comparability, and accumulation of knowledge that empirical research depends upon. It also limits the ability of journalists, policymakers, and educators to understand how digital systems shape public discourse and social life. Together, enclosure, exemption, and erosion represent a structural transformation in the foundations of public knowledge. They challenge assumptions underpinning established research disciplines and threaten the capacity of democratic societies to oversee vital systems.

\section*{Three values}
Data science, broadly understood, overlaps with computer and information science, statistics, economics, operations research, accounting, and geography, and with roles such as data engineer, quantitative researcher, and user researcher. The term is both institutionally influential and still undergoing the kind of professionalization that shaped earlier public interest fields~\cite{passiTrustDataScience2018a,greenDataSciencePolitical2021}. At its most pragmatic level, data science is the identification of patterns and generation of explanations from complex datasets. Despite efforts to portray it as neutral and objective, data science is inescapably political and material~\cite{boydCriticalQuestionsBig2012,vandijckDataficationDataismDataveillance2014}. It is embedded within practices, tools, institutions, and ideologies that shape power relations going into and coming out of data-driven systems, and it depends on infrastructures, standards, and labor that collect, store, maintain, transmit, and compute information~\cite{thomerRelationalDataParadigms2020,acker_DataCraftTheory_2019}. Precisely because of these features, it is a generative site for translating the traditions and accomplishments of established public interest professions to contemporary data-driven systems. Against the pressures of enclosure, exemption, and erosion, public interest data infrastructuring advances a constructive agenda grounded in openness, oversight, and ownership to rebuild the conditions for democratic oversight and accountability~\cite{hutchinsonAccountabilityMachineLearning2021,pasquale_black_2016}.

\paragraph*{Openness}
Openness stands in opposition to enclosure. In a period when platforms withdraw support for APIs, governments sunset datasets, and corporations privatize user-generated content, openness asserts that the ability to observe and understand sociotechnical systems is a public necessity~\cite{Freelon_ComputationalResearchPostAPI_2018,Bruns_APIcalypsesocialmedia_2019}. This commitment involves investing in redundant archives, transparent pipelines, shared community data trusts, and durable public repositories capable of preserving access to public-relevant data even as private incentives shift~\cite{schwartz_archives_2002,acker_DataCraftTheory_2019}. Openness also emphasizes methodological legibility: data provenance, documentation, and reproducible workflows that make it possible for more than experts to evaluate how data is collected, transformed, and interpreted~\cite{thomerRelationalDataParadigms2020,passiTrustDataScience2018a}. By building infrastructures that resist the fencing-off of digital commons, openness helps restore the empirical foundations of public oversight and sustains public interest lineages in journalism and law.

\paragraph*{Oversight}
Oversight confronts exemption by insisting that data-driven systems must be governed with the same rigor, transparency, and public responsibility expected of other critical infrastructures~\cite{pasquale_black_2016,hutchinsonAccountabilityMachineLearning2021}. Platforms and AI developers increasingly claim exceptional status, evading scrutiny by appealing to the proprietary nature of their models or the supposed risks of disclosure~\cite{suzor2018digital,obar2020biggest}. Public interest data infrastructuring advances counter-practices: robust auditing methods, documentation standards, impact assessments, and regulatory frameworks capable of testing and verifying algorithmic claims~\cite{hutchinsonAccountabilityMachineLearning2021,selbst2019fairness}. Drawing on institutional lessons from financial auditing, product safety regulation, and environmental monitoring, oversight treats data-driven systems as subjects of public concern rather than private fiefdoms. In doing so, it establishes conditions under which algorithmic processes can be evaluated, contested, and governed.

\paragraph*{Ownership}
Ownership serves as the counter-force to erosion. As data disappears, infrastructures decay, and fragmented platforms undermine reproducibility, public interest data infrastructuring emphasizes the importance of resilient, public stewardship of data resources~\cite{Freelon_ComputationalResearchPostAPI_2018,gaffney_caveat_2018}. Ownership here does not refer to private possession but to collective governance, in which communities, researchers, public institutions, and civil society organizations participate in establishing norms for long-term data preservation and use~\cite{schwartz_archives_2002,stoler_ColonialArchivesArts_2002}. This involves building interoperable infrastructures, investing in reproducibility frameworks, creating archival strategies that safeguard public memory, and supporting institutions capable of maintaining data resources over decades~\cite{acker_DataCraftTheory_2019,Plantin_Infrastructurestudiesmeet_2016}. By cultivating forms of ownership that distribute responsibility and ensure continuity, public interest data infrastructuring preserves the empirical record essential for democratic accountability and intergenerational justice.

\section*{Public interest data infrastructuring}
The pressures of enclosure, exemption, and erosion do not merely constrain research access; they reshape the institutional conditions under which public knowledge is produced, verified, and contested. Enclosure reorganizes who can observe by redesigning interfaces, archives, and rights-to-know so that scrutiny becomes costly or infeasible in the post-API environment~\cite{Freelon_ComputationalResearchPostAPI_2018,Bruns_APIcalypsesocialmedia_2019,Puschmann_endwildwest_2019}. Exemption reorganizes who must answer by allowing firms and agencies to operate behind legal and technical shields that treat opacity as default and auditing as exceptional~\cite{pasquale_black_2016,Halavais_Overcomingtermsservice_2019}. Erosion reorganizes what can be remembered by degrading archives, standards, and measurement continuity through link rot, platform churn, and the fragility of informal research infrastructures~\cite{gaffney_caveat_2018,baumgartner_incompleteness_2018}. Public interest data infrastructuring responds with an agenda: openness as durable legibility of evidentiary traces; oversight as routinized, testable scrutiny; and ownership as stewardship obligations and collective governance that persist beyond shifting incentives and governance regimes~\cite{vandervlist_APIGovernanceCase_2022}.

\textit{Infrastructuring} is the mechanism that translates these values into durable accountability capacities. Rather than treating infrastructure as a static substrate, infrastructuring emphasizes how accountability is produced through ongoing alignment of standards, tools, routines, and institutions especially in domains where data work is inseparable from organizational practice, maintenance, and trust~\cite{Plantin_Infrastructurestudiesmeet_2016,acker_DataCraftTheory_2019,passiTrustDataScience2018a}. Under conditions of enclosure, exemption, and erosion, the most consequential failures of data-driven governance are infrastructural: access pathways collapse, audit trails are not retained, documentation is missing, and responsibility diffuses across vendors and agencies until harms become untraceable and therefore unaccountable~\cite{passiTrustDataScience2018a}. 

In infrastructure studies, the ``installed base'' names the already-existing socio-technical assemblage—data formats, identifiers, legacy systems, institutional routines, professional jurisdictions, and tacit work practices—that any new information infrastructure must inherit rather than replace wholesale: it is simultaneously a resource (what makes coordination possible) and a constraint (what makes redesign difficult). Star and Ruhleder emphasize that infrastructure is relational and becomes “real” in use, so interventions succeed when they attach to existing practices and breakdown/repair cycles rather than assuming greenfield design; the installed base is therefore the terrain on which access, standards, and accountability are negotiated in practice~\cite{starStepsEcologyInfrastructure1994a,starEthnographyInfrastructure1999a}. Building on this lineage, Aanestad and colleagues foreground “infrastructuring” as ongoing alignment work of incremental extensions, gateways, and governance arrangements that bootstraps new capacities from what is already there, instead of attempting disruptive replacement that institutions cannot sustain~\cite{aanestadInformationInfrastructuresChallenge2017,pipekInfrastructuringIntegratedPerspective2009a}. This perspective also clarifies why public-interest reforms in data systems fail when they stop at principles: values like openness, oversight, and ownership must be translated into durable elements of linkable records, auditable provenance, retention routines, tiered access, and escalation pathways that persist through procurement churn and political turnover~\cite{hansethDesignTheoryDynamic2010,karastiStudyingInfrastructuringEthnographically2018a}. Together, these components define infrastructuring as a practical agenda: building the installed base that makes openness durable, oversight routine, and ownership accountable to publics across generations.

\paragraph{Identifiers.}
Accountability depends on the ability to link evidence across time, agencies, systems, and decision points. Identifiers are persistent, stable references like unique IDs for documents, policies, contracts, model versions, incidents, and actors that enable records to be joined, traced, and cited without ambiguity. Without stable identifiers, archives fragment into unconnectable artifacts: decisions cannot be followed to outcomes, contracts cannot be tracked across amendments, and system changes cannot be audited longitudinally. Infrastructuring identifiers is therefore a standards and governance project as much as a technical one, shaped by classification systems and institutional commitments to durable reference~\cite{bowker_sorting_2000}.

\paragraph{Provenance.}
Openness is not only access to files but access to interpretability. Provenance refers to documentation that makes data and model outputs legible: where records came from, how they were collected, what transformations were applied, what assumptions were embedded, and where uncertainty enters. Provenance enables verification and contestation by turning an output into an inspectable claim rather than an inscrutable result, and it helps sustain trust across heterogeneous data practices~\cite{thomerRelationalDataParadigms2020,passiTrustDataScience2018a}. Infrastructuring provenance means adopting documentation standards, logging transformations, and recording version histories so that downstream users can evaluate what a dataset or model can and cannot support.

\paragraph{Retention.}
Erosion is often mundane: portals disappear, storage media decay, contracts lapse, and formats become unreadable. Retention names the policies and technical mechanisms that preserve the empirical record against loss. Retention includes explicit schedules, redundancy and mirroring, migration plans, and custodial obligations that survive organizational turnover and vendor change. Infrastructuring retention is what makes long-horizon oversight possible, especially when research-relevant datasets depend on fragile, quasi-voluntary infrastructures and shifting platform governance~\cite{acker_DataCraftTheory_2019,gaffney_caveat_2018,baumgartner_incompleteness_2018}.

\paragraph{Access tiers.}
Openness and safety must be reconciled rather than treated as a binary choice. Access tiers are governance arrangements that specify what is public, what is restricted, and under what conditions trusted intermediaries may access sensitive traces. Tiered access supports democratic scrutiny while reducing privacy risks, retaliation risks, and the possibility that transparency becomes a vector for harassment or extraction~\cite{gibney_PrivacyHurdlesThwart_2019,Mervis_Privacyconcernscould_2019}. Infrastructuring access tiers also requires explicit data governance like role-based permissions, documented decision rules, and accountability for intermediaries so that restricted access does not become a pathway for exemption~\cite{khatri_governance_2010}.

\paragraph{Escalation pathways.}
Evidence matters only if it can trigger remedy. Escalation pathways are the interfaces and procedures that connect records to action: complaint channels, audit triggers, appeals and due process mechanisms, public reporting cadences, and enforcement hooks. These pathways ensure that contestation is not merely rhetorical; they provide routes by which affected parties can challenge decisions and by which oversight bodies can compel response. Infrastructuring escalation pathways means making them usable, documented, and routinized, drawing on emerging work in algorithmic accountability and the institutional design of oversight~\cite{hutchinsonAccountabilityMachineLearning2021,pasquale_black_2016}.

\section*{Public interest lineages}
The idea of public interest data infrastructuring does not emerge in a vacuum. Many professions have cultivated explicit public interest missions by developing norms, institutions, and ethical frameworks to ensure that specialized expertise serves the common good rather than private power~\cite{stapletonWhoHasInterest2022,dantecInfrastructuringFormationPublics2013,costanza-chockDesignJustice2020}. Journalism, law, accounting, planning, and engineering each grappled with harms created by opaque systems, unregulated markets, or concentrated authority. Their responses offer important lessons for public interest data infrastructuring, which now faces analogous challenges in the realm of data and algorithms.

\paragraph*{Journalism}
Journalism is rooted in the democratic ideal that an informed public is necessary for self-governance~\cite{ettemaJournalismReasongivingDeliberative2007,deuzeWhatJournalismProfessional2005}. Over time, journalism developed a watchdog role: scrutinizing government, exposing corruption, and making complex information legible to citizens. Legal protections for press freedom and norms such as verification, independence, and accountability emerged to institutionalize these commitments. Investigative reporting made extensive use of public records, freedom-of-information laws, interviews, and data analysis to uncover hidden harms~\cite{glasserInvestigativeJournalismMoral1989,jacquetteJournalismEthicsTruthTelling2009,protessJournalismOutrageInvestigative1992}. Journalism's public interest tradition demonstrates that access to information is a democratic requirement~\cite{zengGhostsNewspapersPublic2022}. For data science, the parallel is clear: without durable rights to observe, analyze, and publish information about digital systems, the public cannot meaningfully evaluate their impacts. Read as infrastructuring, journalism's distinctive contribution is to maintain the evidentiary surface of public life of archives, provenance, verification routines, and publication practices keep systems observable (openness) and contestable (oversight) even when powerful actors prefer opacity.

\paragraph*{Law}
Law developed its own public interest orientation through legal aid movements, civil rights litigation, pro bono service, and cause lawyering~\cite{council1976balancing,rhode2004probono,sarat1998cause,cummings2009public}. Public interest lawyering emerged in response to structural inequality, recognizing that formal rights are meaningless without institutions that help people exercise them~\cite{southworth2005what,cummings2006after,trubek2011public}. Lawyers working in the public interest developed strategic litigation, community partnerships, and impact-driven advocacy to challenge discriminatory policies, regulate dangerous industries, and expand civil liberties~\cite{cummings2009public,nielsen2006organization,albiston2017public}. They also designed institutions such as public defender systems, legal clinics, and nonprofit advocacy organizations that could sustain this work outside market pressures~\cite{rhode2004probono,albiston2017public}. Public interest data infrastructuring must similarly confront problems that are not neutral or evenly distributed: just as public interest lawyers challenge powerful actors, public interest data scientists must illuminate opaque algorithmic systems, quantify disparate impacts, and support communities seeking redress~\cite{barocas2016big,selbst2019fairness}. As infrastructuring, public interest law translates values into enforceable procedures like rights to access, duties to disclose, standards of review, and remedies that make oversight a routinized capability and anchor stewardship obligations over time (ownership).

\paragraph*{Accounting}
Accounting provides one of the clearest historical analogies~\cite{carnegieCriticalInterpretiveHistories1996,leeProfessionalizationAccountancyHistory1995}. In the early twentieth century, recurrent financial crises revealed that private auditing practices routinely failed to protect investors or the public. In response, the profession articulated a explicit "public interest mandate": auditors were to serve not the firms that hired them, but the broader public whose well-being depended on reliable financial information~\cite{quayleWhistleblowingAccountingPublic2021,bakerWhatMeaningPublic2005,dellaportasReflectionsPublicInterest2008}. Regulatory institutions such as the Securities and Exchange Commission (SEC) established standards for disclosure, independence, and verification. Algorithmic systems resemble corporate balance sheets that can hide risks, mask errors, or facilitate fraud unless subject to independent audit. Public interest data infrastructuring can draw directly from accounting's lessons to design methods for algorithmic auditing, documentation, stress-testing, disclosure, and transparency that serve public interests rather than the entities deploying the systems. As an infrastructuring lineage, auditing and accounting formalize the production of legible records, standardized controls, and independent verification that map directly onto algorithmic logging, documentation, and impact-assessment regimes (oversight).

\paragraph*{Planning}
In twentieth-century cities, urban planning practices were often technocratic and exclusionary, producing highways, zoning policies, and redevelopment projects that disproportionately harmed marginalized communities~\cite{jacobs1961death,doryClashUrbanPhilosophies2018,guttenbergPlanningIdeology2009}. The advocacy planning movement emerged in response, insisting that planners had obligations to represent underserved populations, engage impacted communities, and make planning processes more participatory~\cite{dadashpoorDefiningPublicInterest2021}. Planning has also shifted toward more holistic and ecological perspectives that recognize problems such as traffic, affordability, and climate risk as deeply entangled~\cite{hirtPremodernModernPostmodern2009,roheLocalGlobalOne2009}. Planning's public interest tradition highlights that expertise must be paired with democratic responsiveness and an awareness of power~\cite{lanePublicParticipationPlanning2005}. For public interest data infrastructuring, the parallel is the need for participatory design, community stewardship of data, and recognition that technical decisions about what data to collect, what models to build, and how to interpret results have distributive consequences. In infrastructuring terms, planning foregrounds how infrastructures allocate benefits and burdens across space and time, and it supplies participatory governance repertoires for negotiating who owns, maintains, and can contest durable systems (ownership).

\paragraph*{Engineering}
Engineering likewise evolved a public interest ethos, especially in response to infrastructure failures, industrial accidents, and technological risks in the nineteenth and twentieth centuries~\cite{petroski1992engineer,perrow1984normal,vaughan1996challenger,dekkerDriftFailureHunting2016}. Engineering societies developed codes of ethics requiring practitioners to hold paramount the safety, health, and welfare of the public alongside adopting practices of risk assessment~\cite{nathwaniThreePrinciplesManaging1995,boudiaIntroductionRiskRisk2007}. Regulatory bodies emerged to ensure that bridges, buildings, railroads, cars, pharmaceuticals, and medical devices met rigorous safety standards through testing, certification, and oversight of occupational health and safety~\cite{abramsShortHistoryOccupational2001}. Managing complex systems required cultural shifts to embrace practices of safety and high reliability rather than techno-solutionism~\cite{weickOrganizingHighReliability1999,haleWorkingRuleWorkingI2013,haleWorkingRuleWorkingII2013}. Public interest data infrastructuring should emulate engineering's commitments to safety, standards, and oversight: data and algorithmic systems increasingly function as critical infrastructure and should be governed accordingly. As infrastructuring, safety engineering emphasizes standards, testing, incident learning, and high-reliability operations as capabilities that operationalize oversight under complexity and create feedback loops that keep infrastructures governable rather than merely optimized.

\section*{Public interest conversations}
Public interest data infrastructuring shares intellectual ground with several adjacent fields but is distinct in its commitments, methods, and institutional ambitions. These neighboring traditions help frame important ethical, technical, and societal concerns around data systems, but they often approach these concerns with goals other than building durable infrastructures and regulatory regimes for public accountability. This section reads adjacent conversations through an infrastructuring lens: what durable capacities do they build (or fail to build) for keeping data- and model-mediated systems observable, auditable, and collectively governed in the face of enclosure, exemption, and erosion?

\paragraph*{Public interest technology}
Public Interest Technology (PIT) is one of the most important of these adjacent fields. PIT refers to a set of practices for designing, deploying, and governing technology in ways that advance the public interest, with explicit commitments to human rights, fairness, and social well-being~\cite{DefiningPublicInterest,abbasSocioTechnicalDesignPublic2021,washingtonDefiningPublicInterest2024}. It is inherently interdisciplinary, drawing together technologists, policymakers, activists, and affected communities to grapple with the ethical, legal, policy, social, economic, and political implications of technology~\cite{abbasCoDesigningFuturePublic2021,chowdharyCareLiberationCreating2020,stapletonWhoHasInterest2022,chambersBigTechAdvocacy2025}. Historically, PIT has been built through collaborations among civil society organizations, foundations, and universities~\cite{kramerBuildingFieldPublic,luskPublicInterestTechnology2021}. Philanthropic initiatives such as the Public Interest Technology University Network (PIT-UN) have invested in fellowships, curricular experiments, and talent pipelines aimed at producing ``civic-minded technologists and digitally fluent policy leaders''~\cite{luskPublicInterestTechnology2021,kramerBuildingFieldPublic}. Whereas PIT takes technology in general as its object and concentrates on building career pathways, capacity, and ethical sensibilities across sectors~\cite{mcguinnessPowerPublicPromise2021,knodelShapingProfessionBuilding2025}, public interest data infrastructuring focuses more specifically on data and algorithmic systems as critical infrastructures that mediate public authority and public accountability~\cite{alkhatibStreetLevelAlgorithmsTheory2019,passiTrustDataScience2018a,hutchinsonAccountabilityMachineLearning2021}. It centers questions of measurement, archiving, auditing, and governance: how to construct evidence commons, dataset and labeling accountability practices, and oversight mechanisms that remain legible to institutions and contestable by publics~\cite{mullerDesigningGroundTruth2021,paakkonenBureaucracyLensAnalyzing2020b,greenAlgorithmicRealismExpanding2020}. In that sense, PIT supplies much of the professional, pedagogical, and coalition-building scaffolding, while public interest data infrastructuring specifies a more targeted agenda for governing data and models as part of the democratic commons~\cite{stapletonWhoHasInterest2022,mcguinnessPowerPublicPromise2021,hutchinsonAccountabilityMachineLearning2021}. From an infrastructuring perspective, PIT contributes coalition-building and professional pathways that make long-horizon stewardship work possible; its limitation is that capacity-building without enforceable recordkeeping and oversight can leave accountability dependent on goodwill.

\paragraph*{Data for good}
Data and AI for social good initiatives (often branded as ``data for good'') have largely grown through philanthropic, corporate, and nonprofit programs that apply analytics to public problems like poverty, education, public health, and disaster response via time-bounded projects and cross-sector partnerships~\cite{williamsDataActionUsing2022,cowlsAISocialGood2021,aulaSteppingBackData2023}. These efforts can deliver rapid analytical capacity to organizations with limited technical resources and help legitimate evaluation criteria oriented toward social benefit rather than solely organizational performance~\cite{cowlsAISocialGood2021}. Yet this same orientation can make the work tool-centric and episodic: success is defined by deploying models, dashboards, or decision-support systems that help institutions ``manage'' a problem, while leaving the underlying political economy of data collection, representation, and harm largely intact~\cite{floridiHowDesignAI2020}. Moreover, many programs implicitly rely on continuing goodwill from platforms and private vendors for access to data and infrastructure, an assumption destabilized by the ``post-API'' environment and by shifting arrangements between public agencies and commercial intermediaries~\cite{Freelon_ComputationalResearchPostAPI_2018,Bruns_APIcalypsesocialmedia_2019,taylorPublicActorsPublic2021}. Public interest data infrastructuring diverges by treating socially beneficial applications as insufficient without durable accountability infrastructures: maintained administrative records, auditable pipelines, public reporting regimes, and oversight institutions designed to resist enclosure, exemption, and selective visibility over time~\cite{schwartz_archives_2002,stoler_ColonialArchivesArts_2002,acker_DataCraftTheory_2019}. Through an infrastructuring lens, ``data for good'' is strongest at rapid problem-solving but weakest at sustaining the installed base like archives, standards, and governance needed for accountability after partners, interfaces, or incentives change.

\paragraph*{Digital government}
``GovTech'' is often framed as the modernization of public administration through digital service delivery, open data infrastructures, participatory design, and platform logics, drawing on broader accounts of public-sector digital transformation, service-maturity models, and platform-oriented reform agendas~\cite{eomDigitalGovernmentTransformation2022,layneDevelopingFullyFunctional2001,brownAppraisingImpactRole2017}. These shifts exist alongside concerns about whether technological change actually strengthens transparency and accountability or instead re-routes them through new intermediaries, opaque integrations, and vendor-mediated infrastructures~\cite{bertotUsingICTsCreate2010,bharosaRiseGovTechTrojan2022,silvePoliticalEconomyGovTech2023}. A public interest lens makes this tension explicit: digital government projects are not only engineering endeavors, but institutional redesigns that decide who is seen, what is measured, and how contestation is handled across complex, ``wicked'' administrative systems~\cite{gil-garciaDigitalGovernmentPublic2018,kimDigitalGovernmentWicked2016,twizeyimanaPublicValueEGovernment2019}. Triangulating with public interest data infrastructuring clarifies the operational requirements for durable accountability infrastructures like auditable administrative records, documentation practices, public reporting regimes, and independent oversight that persists beyond launch and procurement cycles~\cite{panagiotopoulosPublicValueCreation2019,zuiderwijkOpenDataPolicies2014b,bertotUsingICTsCreate2010}. It also requires confronting the political economy of GovTech around procurement, vendor lock-in, and limited competition, dynamics that can quietly shift governing power from public institutions to contractors while reducing accountability to affected communities~\cite{silvePoliticalEconomyGovTech2023,bharosaRiseGovTechTrojan2022,brownAppraisingImpactRole2017}. GovTech sits closest to infrastructuring because it operates where procurement, standards, and administrative records are made; the central question is whether modernization builds public auditability and stewardship or deepens vendor-mediated exemption.

\paragraph*{Civic technology}
Civic technology shares a complementary but distinct lineage, focused on building and maintaining digital tools that improve public service delivery, enable citizen engagement, and streamline administrative processes across government and civil society~\cite{boehnerDataDesignCivics2016,aragonCivicTechnologiesResearch2020,ruijerConnectingSocietalIssues2017}. Civic tech has also contributed to modernization efforts by prototyping new participation modalities and helping consolidate institutional and professional identities ranging from ``civic technologist'' roles inside government to hybrid communities of practice that blend design, engineering, and public problem-solving~\cite{aragonCivicTechnologiesResearch2020,matias_civic_2019,riderBuildingIdealWorkplaces2022}. Yet civic technology can privilege efficiency and near-term delivery over deeper questions of fairness, data governance, and algorithmic accountability, especially when complex sociotechnical systems are abstracted into narrow problem definitions and success metrics~\cite{selbst2019fairness,hutchinsonAccountabilityMachineLearning2021,paakkonenBureaucracyLensAnalyzing2020b}. It also tends to operate within existing institutional arrangements and their procurement constraints, legacy information systems, and established bureaucratic mandates rather than contesting or redesigning them, which can limit its ability to shift underlying power relations even when participation interfaces improve~\cite{kimDigitalGovernmentWicked2016,gil-garciaDigitalGovernmentPublic2018,ruijerConnectingSocietalIssues2017}. Public interest data infrastructuring, treats civic capacity not only as a matter of better tools, but as a matter of building and reforming institutions for record-keeping, archiving, auditing, and stewardship so that public knowledge remains accessible, contestable, and durable over time~\cite{acker_DataCraftTheory_2019,schwartz_archives_2002,hutchinsonAccountabilityMachineLearning2021}. Civic tech often excels at participatory interfaces and delivery prototypes, but infrastructuring asks what remains after the pilot: who maintains the data, what standards persist, and what oversight pathways are institutionalized.


\paragraph*{Critical data studies}
Critical data studies offers an essential set of theoretical and interpretive tools for analyzing how data practices construct categories of identity, expand surveillance capacities, and reproduce or reconfigure power relations in everyday life and institutional governance~\cite{boydCriticalQuestionsBig2012,vandijckDataficationDataismDataveillance2014,puschmannBigDataBig2014}. Scholars in this tradition illuminate the political, historical, and ethical stakes of datafication by showing how ``raw'' data are socially produced, how classificatory regimes can enact harm, and how algorithmic systems embed and operationalize inequality through design choices, training data, and institutional contexts~\cite{dentonGenealogyMachineLearning2021,greenDataSciencePolitical2021}. Public interest data infrastructuring stands on these shoulders, particularly in recognizing how enclosure, exemption, and erosion in contemporary data ecosystems threaten the durability and contestability of democratic knowledge~\cite{Freelon_ComputationalResearchPostAPI_2018,taylorPublicActorsPublic2021,Bruns_APIcalypsesocialmedia_2019}. Yet critical data studies is often organized around diagnosis and critique, and is less consistently oriented toward specifying implementable institutional, infrastructural, and regulatory alternatives that can be sustained within public governance~\cite{latourWhyHasCritique2004,neffCritiqueContributePracticeBased2017}. Public interest data infrastructuring extends the critical tradition by translating analysis into operational commitments of auditable records, documented pipelines, stewardship regimes, and oversight mechanisms grounded in precedents from other public interest professions and designed to hold up under changing business models and political winds~\cite{acker_DataCraftTheory_2019,schwartz_archives_2002,hutchinsonAccountabilityMachineLearning2021}. Critical data studies supplies indispensable diagnostic language for infrastructuring to ask how power shapes classification, surveillance, and datafication but it is most complementary when paired with institutional designs that translate critique into durable accountability machinery.

\paragraph*{Data justice}
Data justice treats datafication as a social and political project rather than a neutral technical evolution, centering questions of power, redistribution, and representation in data relations~\cite{taylorWhatDataJustice2017,dignazioDataFeminism2020,braunDataJusticeData2022}. It also emphasizes lived experience, historical inequalities, and communities' capacity to refuse, contest, or re-shape harmful systems and extractive data regimes~\cite{zongDataRefusalFramework2024,garciaNoReimaginingData2022,dencikRescuingDataJustice2025}. Because data justice scholarship and advocacy can at times operate at some distance from the day-to-day design, procurement, and maintenance of statistical and administrative infrastructures, public interest data infrastructuring names a complementary project: it shares data justice's normative commitments while focusing on building and governing concrete infrastructures that embed those commitments in institutional routines~\cite{taylorWhatDataJustice2017,acker_DataCraftTheory_2019,hutchinsonAccountabilityMachineLearning2021}. In this division of labor, data justice contributes critical lenses and movement-led agendas that define what counts as the public interest~\cite{dignazioDataFeminism2020,braunDataJusticeData2022}, while public interest data infrastructuring develops methods and institutional footholds of auditable records, maintained archives, documentation standards, and oversight mechanisms needed to contest unjust datafication and sustain democratic accountability at scale~\cite{schwartz_archives_2002,acker_DataCraftTheory_2019,hutchinsonAccountabilityMachineLearning2021}. Data justice specifies what infrastructuring must be for: redistribution, representation, and the capacity to refuse; public interest data infrastructuring builds the records, governance, and oversight arrangements that can realize these aims at scale.

\section*{Anticipating critiques}

Any agenda that tries to reshape the political economy of data will encounter resistance. Public interest data infrastructuring challenges familiar arrangements of data access, corporate secrecy, regulatory design, and epistemic authority. The point of anticipating critiques is not simply rhetorical defense; it is to treat objections as \emph{design constraints} on the installed base. Each critique identifies a different failure mode like missing identifiers, absent provenance, weak retention, unsafe access, or ineffective escalation and thereby clarifies what infrastructuring must build to make openness, oversight, and ownership durable.

\subsection*{Realist}
\begin{quote}\centering
\textit{``The open data era is over and never coming back.''}
\end{quote}

\noindent The realist objection begins from an uncomfortable truth: platforms that once courted developers and researchers with generous APIs now treat data as a strategic asset, and governments have withdrawn or underfunded key datasets. From this vantage point, calls for renewed openness can look nostalgic; enclosure appears as the natural maturation of technological and economic systems. Public interest data infrastructuring accepts the descriptive force of this story but rejects its fatalism. Historically, sectors far more opaque and powerful than today's platforms have been opened to scrutiny not by goodwill, but by publics, journalists, lawyers, and regulators building new accountability infrastructures.

Realism relocates the work from voluntary access to the installed base. Under enclosure and erosion, openness becomes a retention and reference problem: build \emph{redundant repositories} and enforce \emph{retention} so evidence does not disappear; require \emph{persistent identifiers} and stable citation surfaces so public-relevant traces can be linked even as interfaces change; and treat \emph{provenance} as mandatory so that externally visible data can be interpreted and contested. Where full openness is infeasible, tiered \emph{access} with accountable intermediaries can create legitimate oversight channels without restoring the old API world. Finally, realism requires \emph{escalation pathways}: statutory access rights, mandated research portals, and audit triggers that convert noncooperation into enforceable obligations rather than a dead end. Realism is valuable when it reveals the limits of corporate goodwill; it becomes problematic only when it slides into resignation.

\subsection*{Critical}
\begin{quote}\centering
\textit{``How is this different from existing reformist projects?''}
\end{quote}

\noindent From data justice, critical data studies, and long-standing public interest technology communities comes proliferation fatigue: ``public interest data infrastructuring'' can sound like just another label layered atop data for good, civic tech, and responsible AI. This critique matters because it forces clarity about distinctiveness. The claim is not that public interest data infrastructuring is wholly new, but that it integrates and redirects existing strands toward institutional change and long-horizon maintenance.

The distinctiveness of public interest data infrastructuring is not a new moral vocabulary but a design stance: treat accountability as an installed base that must be built, governed, and maintained. This stance is testable. If a project cannot specify identifiers, provenance, retention, access tiers, and escalation pathways, it is likely another pilot that cannot survive enclosure and erosion. The critical critique also sharpens \emph{whose} installed base is being built: provenance must include the social assumptions and category choices embedded in measurement; access tiers must avoid turning openness into surveillance; and escalation pathways must include routes for affected communities to contest harms. In this sense, critical perspectives are not external to infrastructuring but internal guardrails: they define what counts as public interest and ensure that infrastructuring does not reproduce extractive visibility or technocratic closure under the banner of accountability.

\subsection*{Libertarian}
\begin{quote}\centering
\textit{``Oversight will kill innovation.''}
\end{quote}

\noindent The libertarian critique reprises a familiar argument: tighter oversight and disclosure will slow innovation, deter investment, and erode competitiveness. In data and AI, this appears as warnings that audit mandates, transparency rules, or access obligations will chill experimentation and burden startups more than incumbents. Poorly crafted regimes can indeed entrench large firms and overwhelm smaller actors. But the broader claim that accountability and innovation are inherently at odds does not withstand historical scrutiny across safety, finance, and environmental governance.

Infrastructuring reframes ``oversight'' as a minimal set of accountability affordances that often \emph{reduce} friction by standardizing expectations. Rather than mandating full disclosure of proprietary systems, it emphasizes \emph{provenance} and \emph{logging} sufficient to verify claims, \emph{identifiers} and versioning to track changes, and standardized \emph{interfaces for audit} that lower the transaction costs of scrutiny for both regulators and firms. Retention requirements can be scoped to what is necessary for incident reconstruction, while access tiers protect sensitive information without creating blanket exemption. Most importantly, infrastructured oversight creates predictable escalation pathways like incident reporting, audit triggers, and appeal mechanisms that can prevent catastrophic failures and the blunt political backlash they invite. The libertarian critique therefore becomes a design challenge: build lightweight, interoperable, and iterative installed base requirements that enable competition on robustness rather than secrecy.

\subsection*{Positivist}
\begin{quote}\centering
\textit{``Data science should not be political.''}
\end{quote}

\noindent Technical disciplines often voice concerns that explicitly orienting data science toward justice-centered values risks politicizing a field that should remain method-driven. In this positivist critique, algorithms are tools; ethics should reside in downstream use, not in methods. The worry is that treating every modeling decision as political invites relativism and paralyzes practice. Public interest data infrastructuring treats this critique as a demand for operational clarity: what, precisely, would make claims about data-driven systems more testable?

Infrastructuring aligns with positivist commitments by strengthening the conditions for independent verification. The argument is not that every choice is partisan, but that pipelines are never interpretation-free and that rigor improves when assumptions are documented and outputs are reproducible. Concretely, provenance and documentation make modeling claims inspectable; identifiers and versioning make comparisons meaningful across time; retention sustains the empirical record needed to detect drift; and tiered access enables qualified scrutiny without collapsing privacy into secrecy. Where value disputes arise, escalation pathways channel disagreement into structured review like audits, impact assessments, and appeals rather than informal debate. Infrastructuring thus shifts ``politics'' into a procedural register: it makes claims contestable through evidence, standards, and repeatable methods, rather than relying on authority or trust alone.

\subsection*{Pessimist}
\begin{quote}\centering
\textit{``These systems are too complex to oversee.''}
\end{quote}

\noindent The pessimist critique centers on the opacity of modern data systems. Deep learning models with billions of parameters, constantly updated recommender engines, and sprawling pipelines are difficult to understand even for their creators. If systems are complex and rapidly evolving, skeptics ask, how can auditors, regulators, or the public meaningfully oversee them? Public interest data infrastructuring reframes this as a design and institutional problem, not a conversation stopper. Complex systems in aviation, finance, and environmental governance are overseen through layered standards, specialized roles, structured disclosure, and institutional memory.

Complexity becomes governable when the installed base supports \emph{auditable complexity}. Identifiers and change logs allow model and system evolution to be tracked; provenance documents data lineage, training regimes, and key assumptions; retention ensures incidents can be reconstructed and longitudinal impacts evaluated. Access tiers enable role-based visibility, so independent intermediaries can examine sensitive traces while remaining accountable to the public. Escalation pathways routinize response like incident reporting, stress testing, postmortems, and appeals so oversight is continuous rather than episodic. The pessimist critique is therefore a demand for infrastructuring discipline: build feedback loops and institutional memory that make scrutiny possible even when full transparency or full interpretability is unattainable. Complexity is precisely why infrastructuring matters: without maintained records and routinized escalation, ``too complex to oversee'' becomes a self-fulfilling exemption.

% \section*{Lineages and conversations against critiques}

% Table~\ref{tab:lineages-critiques-confusion} summarizes how public interest lineages and adjacent conversations (rows) respond to anticipated critiques (columns). In this revision, we read each row not only as a normative stance or methodological repertoire, but as a partial account of \emph{infrastructuring}: the work of building and maintaining an installed base that makes data-driven systems observable, auditable, and contestable over time. We therefore interpret each critique as a demand on specific installed base components. Realist critiques stress whether accountability can persist under access loss, fragmentation, and institutional decay: retention, identifiers, and resilient access tiers. Critical critiques stress whose interests the installed base serves and how it distributes burdens: participatory governance of provenance, access tiers, and community-centered escalation pathways. Libertarian critiques stress whether oversight becomes overbroad constraint: minimal, standardized provenance and audit interfaces plus narrowly scoped retention and access tiers. Positivist critiques stress testability and methodological clarity: provenance, identifiers, and reproducible retention practices. Pessimist critiques stress governability under complexity: layered access tiers, routinized escalation pathways, and institutional memory grounded in retention and versioned identifiers.

% \begin{table*}[t]
% \centering
% \small
% \setlength{\tabcolsep}{6pt}
% \renewcommand{\arraystretch}{1.25}
% \begin{tabular}{p{1.9in}ccccc}
% \toprule
% \textbf{Lineage + conversation} & \textbf{Realist} & \textbf{Critical} & \textbf{Libertarian} & \textbf{Positivist} & \textbf{Pessimist} \\
% \midrule
% Journalism (watchdog, FOI, verification) & $\uparrow\uparrow$ & $\downarrow$ & $\uparrow$ & $\uparrow\uparrow\uparrow$ & $\uparrow$ \\
% Law (legal aid, rights, litigation, due process) & $\uparrow\uparrow$ & $\uparrow$ & $\uparrow\uparrow$ & $\uparrow\uparrow\uparrow$ & $\uparrow\uparrow$ \\
% Accounting (audit mandate, disclosure regimes) & $\uparrow\uparrow\uparrow$ & $\uparrow$ & $\uparrow\uparrow\uparrow$ & $\uparrow\uparrow$ & $\uparrow\uparrow$ \\
% Planning (advocacy planning, participation, power) & $\uparrow\uparrow$ & $\downarrow$ & $\uparrow$ & $\uparrow\uparrow\uparrow$ & $\uparrow$ \\
% Engineering (safety culture, standards, high reliability) & $\uparrow\uparrow\uparrow$ & $\uparrow$ & $\uparrow\uparrow\uparrow$ & $\uparrow\uparrow$ & $\uparrow\uparrow\uparrow$ \\
% \midrule
% Public interest technology (field-building, ethics, pathways) & $\uparrow\uparrow$ & $\uparrow\uparrow$ & $\uparrow$ & $\uparrow\uparrow\uparrow$ & $\uparrow$ \\
% Data for good (projects, partnerships, applications) & $\uparrow$ & $\downarrow\downarrow$ & $\downarrow$ & $\downarrow\downarrow$ & $\downarrow$ \\
% Govtech (service reform, platforms, procurement) & $\uparrow\uparrow$ & $\uparrow$ & $\uparrow\uparrow$ & $\downarrow$ & $\uparrow\uparrow$ \\
% Civic technology (tools, participation interfaces, delivery) & $\uparrow$ & $\uparrow$ & $\uparrow\uparrow$ & $\downarrow$ & $\uparrow$ \\
% Critical data studies (power, surveillance, datafication critique) & $\uparrow\uparrow$ & $\uparrow\uparrow\uparrow$ & $\downarrow\downarrow$ & $\uparrow\uparrow\uparrow$ & $\downarrow\downarrow$ \\
% Data justice (redistribution, representation, refusal, lived experience) & $\uparrow\uparrow$ & $\uparrow\uparrow\uparrow$ & $\downarrow\downarrow$ & $\uparrow\uparrow\uparrow$ & $\downarrow\downarrow$ \\
% \bottomrule
% \end{tabular}
% \caption{A confusion matrix of how public interest lineages and adjacent conversations (rows) perform against anticipated critiques of a public interest data infrastructuring agenda (columns). Cell values indicate relative strength/weakness in addressing each critique: \textbf{$\uparrow\uparrow\uparrow$} very strong to \textbf{$\downarrow\downarrow\downarrow$} very weak.}
% \label{tab:lineages-critiques-confusion}
% \end{table*}

% To populate the matrix, we interpret each critique as an evaluative demand on infrastructuring capacity: whether a lineage's ``default toolkit'' can specify, build, and govern the installed base required for durable accountability. For each row, we summarized its characteristic repertoire (\textit{e.g.}, journalism's verification and FOI routines; accounting's disclosure and audit mandates; data justice's redistribution and refusal). For each column, we operationalized the critique as a set of questions about installed base components (\textit{e.g.}, does this toolkit create retention and identifiers that survive access loss? does it provide escalation pathways and governance over access tiers? does it strengthen provenance and versioning for testability?). Each cell reflects an abductive assessment of fit: not an idealized best case, but the extent to which the lineage, \emph{on its own terms}, tends to furnish the infrastructural prerequisites of openness, oversight, and ownership.

% Three grounded clusters remain visible, but the installed base lens clarifies what each cluster contributes and where coalitions are structurally necessary. No single lineage resolves the critique space because no single tradition spans the full installed base. The most resilient response is a coalition: contestation traditions clarify what the installed base must be \emph{for}, accountability-infrastructure traditions supply enforceable recordkeeping and escalation machinery, and delivery traditions implement interoperable systems that people can actually use. Public interest data infrastructuring is the integrative program that aligns these repertoires so that openness, oversight, and ownership become durable civic capacities

% \paragraph{Accountability infrastructure cluster: law, accounting, engineering.}
% Accounting, engineering, and law are consistently strong against the Realist and Pessimist critiques because they presume governability under adversarial incentives and complexity, and they provide institutional templates for making oversight durable. Read as infrastructuring traditions, they emphasize \emph{identifiers} (standardized reference and versioning), \emph{provenance} (documentation and disclosure), \emph{retention} (recordkeeping obligations), and \emph{escalation pathways} (appeals, enforcement hooks, incident reporting). These are precisely the installed base components most threatened by exemption and erosion, and they map to oversight and ownership as maintained capacities rather than aspirational ideals. The predictable vulnerability of this cluster is that it triggers Libertarian objections because it normalizes constraint; the installed base lens helps refine the response: focus on minimal, interoperable requirements (versioned identifiers, scoped logs, and audit interfaces) rather than maximal disclosure.

% \paragraph{Sensemaking and contestation cluster: critical data studies and data justice.}
% Critical data studies and data justice are strongest against the Critical critique because they define what counts as public interest and keep power, harm, and refusal legible. Through the installed base lens, their distinctive contribution is governance over \emph{provenance} and \emph{access tiers}: they insist that categories, documentation, and visibility regimes distribute burdens and benefits, and that infrastructuring must include participatory decision-making about what is recorded, who can see it, and to what ends. They also strengthen escalation pathways by foregrounding whose claims are recognized and what remedies are legitimate. Their recurring weakness against the Pessimist critique is not conceptual but procedural: on their own, they can under-specify durable apparatuses like retention obligations, standardized identifiers, and routinized audit interfaces that survive procurement cycles and institutional churn. This weakness is precisely where coalition with accountability-infrastructure traditions becomes necessary.

% \paragraph{Delivery and reform cluster: GovTech and civic tech.}
% GovTech and civic tech score well under Realist conditions because they operate at the systems layer where installed base choices are made: procurement, architectures, cross-agency integration, and service delivery. Read as infrastructuring traditions, they are most relevant to \emph{identifiers} (interoperability standards), \emph{retention} (records management in administrative systems), and the design of \emph{access tiers} (portals, dashboards, and mediated disclosure). GovTech tends to outperform civic tech on the Pessimist critique because it engages the administrative backbone rather than only interface-level participation. Their weakness under the Critical critique is also infrastructural: modernization can inadvertently deepen enclosure (via vendor lock-in), strengthen exemption (via proprietary platforms), or widen surveillance (via increased visibility without governance). The installed base framing clarifies the remedy: procurement and platform reforms must explicitly secure provenance, exportability, and public control over indexing and access rules.

% \paragraph{Connective tissue: journalism and public interest technology.}
% Journalism performs strongly for Positivist audiences because verification routines translate normative commitments into publicly legible evidence; in installed base terms, it excels at building and interpreting provenance and at activating escalation through publication and agenda-setting. Its Realist vulnerability is likewise infrastructural: without retention, identifiers, and access rights, verification becomes prohibitively costly. Public interest technology performs well against the Critical critique when it foregrounds ethics, coalition-building, and professional pathways; its installed base contribution is indirect but essential: PIT can supply the institutional capacity and cross-sector coordination required to staff and sustain stewardship, intermediary access tiers, and audit organizations. Where PIT is not coupled to enforceable retention, disclosure, and escalation pathways, it risks becoming a pipeline-and-capacity project without durable infrastructuring impact.

\section*{Three Cases}

Public interest data infrastructuring is not solely a methodological shift but a civic and institutional project to rebuild and maintain the conditions under which democratic societies can understand, evaluate, and govern powerful data-driven systems. Three speculative case studies illustrate how this infrastructuring could play out: local government transparency, platform-mediated work, and climate-relevant memory. Across the cases, we treat infrastructuring as the practical substrate of public interest data science: the ongoing work of building and maintaining an installed base of records, standards, interfaces, and governance routines that make data-driven power observable, auditable, and contestable over time. To keep the analysis comparable, each case is structured as a stress test of the same framework: how enclosure, exemption, and erosion manifest in a domain, and how openness, oversight, and ownership can be operationalized as durable capacities rather than one-off disclosures or tools, and an installed base checklist to evaluate what counts as ``infrastructure'' in each setting.

\subsection*{Case 1: Mining municipal archives}

Local governments sit at the center of a long public interest tradition: open-records laws, public budgeting, zoning decisions, environmental monitoring, police oversight, school governance, and the routine administration of services. Each meeting, permit, inspection, and vote leaves a trail: council packets, land-use maps, police logs, budget ledgers, environmental reports, procurement contracts, housing documents, and voting records. Taken together, these materials form a vast, largely unstructured corpus of human-generated data about how power operates close to home~\cite{barariLocalViewDatabasePublic2023,sumnerCrowdsourcingReliableLocal2020}. On paper, sunshine laws, public records statutes, and open-meetings requirements were designed to ensure that residents, journalists, researchers, and advocacy groups can see how decisions are made and resources allocated~\cite{einsteinWhoParticipatesLocal2019,yoderDoesPropertyOwnership2020,holmanTablingDebateHow2025}. 

In practice, however, the material reality of disclosure often frustrates that promise. Documents are posted as scanned PDFs, scattered across departmental websites, and mediated through brittle portals that make search and reuse difficult at scale~\cite{brownCouncilDataProject2021,barariLocalViewDatabasePublic2023,fuNaturalLanguageProcessing2024}. Audio and video recordings are increasingly available (often through third-party hosting), but remain hard to systematically analyze without substantial extraction and structuring work~\cite{barariPromiseTextAudio2025a,xuPUBLICSPEAKHearingPublic2025}. Other materials are available only through formal requests that can be slow, iterative, and costly.\cite{warrenRequestAtlasSupportingSlow2025} The result is an archive that is simultaneously public and opaque. The records exist, but they might as well be locked in a basement filing cabinet. This case asks what it would take to turn inert, fragmented disclosures into a living evidence base for oversight and public learning, and how emerging computational tools complicate and expand that project~\cite{spangherTrackingNewsworthinessPublic2024,fuNaturalLanguageProcessing2024}.

\subsubsection*{Pressures and values}
Municipal archives reveal enclosure as ``soft enclosure'': records are nominally public but practically fenced off by scanned PDFs, inconsistent metadata, fragile portals, and weak search~\cite{brownCouncilDataProject2021,fuNaturalLanguageProcessing2024}. Openness as infrastructuring implies accessibility and standardization: machine-readable formats, stable identifiers, consistent schemas, and retrieval systems that make discovery and reuse routine rather than heroic. Exemption appears when disclosure is treated as compliance: meeting packets are posted but cannot be traced across policy lifecycles, contracts, budgets, and outcomes. Oversight requires audit-ready recordkeeping to link decisions to money, vendors, and performance so accountability is testable rather than performative. Erosion is constant: link rot, migrations, vendor changes, and quiet disappearance of attachments and historical pages. Ownership counters erosion through stewardship: redundant mirrors, retention policies, professional custodianship by archivists, and governance that maintains continuity across administrations. In this case, infrastructuring ties openness, oversight, and ownership to the installed base that makes ``transparency'' operational.

\subsubsection*{Anticipated critiques}
Capacity is the first critique: many jurisdictions lack staff, budgets, and coherent records practices; proposals for metadata standards or AI-enhanced retrieval can look like urban techno-solutionism. The infrastructuring implication is that ``better disclosure'' must be delivered through shared services, reusable toolchains, and custodial partnerships that reduce marginal effort for small jurisdictions~\cite{warrenRequestAtlasSupportingSlow2025}. Representation and privacy are the second critique: archives reflect what governments choose to document, and better search can increase exposure for vulnerable residents while amplifying enforcement-oriented records. This critique signals that access tiers and redaction workflows should be treated as core installed base components, not afterthoughts, and that communities need governance leverage over visibility and use. Vendor capture is the third critique: commercial indexing can create a new enclosure around public memory and produce ``transparency theater.'' The infrastructuring response is public or community control over indexing and retrieval, transparent ranking and provenance, exportable indices, and escalation pathways that connect improved usability to real oversight rather than aesthetic dashboards.

\subsubsection*{Lineages and conversations}
Several precedents supply the most directly relevant infrastructuring repertoire. Journalism contributes verification routines, FOI practice, and public-facing translation as capabilities that turn sprawling records into contestable evidence and that expose ``soft enclosure'' when records are formally open but practically unusable. Law supplies enforceable access rights, procedural remedies, and due-process constraints that convert archival legibility into oversight leverage while disciplining privacy and exposure risks. Planning contributes participatory governance repertoires for deciding what should be documented, how categories allocate burdens and benefits, and how communities should govern visibility in local records. Finally, GovTech provides the most immediate institutional footholds for procurement rules, standards, retention policies, and cross-agency interoperability where the installed base can be rebuilt at scale rather than through isolated projects.

\subsubsection*{Infrastructuring synthesis}
Public interest data infrastructuring for municipal archives begins by treating local transparency as a lived civic service rather than a legal formality. The goal is not simply to post records, but to make the everyday operations of government \textit{followable}: decisions can be traced, money can be tracked, and outcomes can be compared across time. In practice, this means building a ``public memory layer'' that sits alongside existing city workflows and gradually repairs the soft enclosure created by PDFs, fragmented portals, and brittle search`\cite{brownCouncilDataProject2021}. Rather than asking every resident or journalist to reinvent retrieval, the infrastructure would provide a stable reference surface for the civic record so a council vote, a permit, a contract amendment, and a budget line item can be connected without bespoke detective work~\cite{barariLocalViewDatabasePublic2023}.

Crucially, this is not a purely technical project; it is an institutional one. Stewardship can be organized through libraries, inter-municipal consortia, or university-community partnerships that provide shared services to jurisdictions that cannot sustain sophisticated archives alone. Openness is implemented as usable access like clear documentation, consistent schemas, and retrieval systems that make discovery routine while also acknowledging that local records include sensitive traces. A public interest approach therefore couples improved accessibility with governed visibility: redaction and tiered access are part of the infrastructure, not afterthoughts, and communities have a say in which forms of visibility increase safety and which amplify harm. Finally, AI becomes valuable only as a public tool inside this pipeline: extraction and summarization are constrained by transparent provenance and exportable indices that prevent vendor capture. The success metric is democratic capacity: whether oversight bodies, residents, and journalists can repeatedly substantiate claims, challenge decisions, and hold local power to account without relying on heroic effort or proprietary intermediaries.


\subsection*{Case 2: Supporting worker observatories}

The platformization of work has reconfigured the basic terms of the labor contract. In ride-hailing, delivery, care work, and an expanding array of gig sectors, the most consequential information about workers' lives is collected, interpreted, and acted upon almost entirely by companies. Apps track GPS traces, acceptance rates, completion times, customer ratings, and opaque ``quality'' scores that determine access to shifts, bonuses, and even continued employment~\cite{jarrahiAlgorithmicManagementWork2021}. At the same time, platforms themselves are subject to ``enshittification'': gradual degradation of conditions as firms optimize for investor returns, visible in shifting pay and pricing regimes and other unilateral changes to work conditions~\cite{binnsNotEvenNice2025}. Data becomes both the means and the medium of control. Workers are constantly monitored, sorted, and evaluated, yet have almost no visibility into the metrics that govern them and few avenues for appeal~\cite{jarrahiAlgorithmicManagementWork2021,steinYouAreYou2023}. The political economy of gig work is defined by extreme asymmetries of information and power. Platforms invest heavily in sociomaterial infrastructures of data centers, telemetry pipelines, dashboards, legal teams that convert individual acts of labor into real-time behavioral streams. These infrastructures are tightly coupled to contractual regimes that classify workers as independent contractors, limiting obligations for benefits, job security, or due process. Deactivations, shadow bans, or re-ranking are frequently triggered by automated systems whose logic is treated as proprietary. Appeals, when available, are routed through channels that offer little transparency or recourse.

Against this backdrop, ``worker observatories'' have begun to emerge. Drawing on long traditions of workers' inquiry, mutual aid, and public interest organizing, gig workers and their allies are building grassroots data infrastructures designed to monitor, document, analyze, and resist platform behavior from the bottom up~\cite{calacciOrganizingEndEmployment2022,calacciBargainingBlackBoxDesigning2022,vincentDataStrikesEvaluating2019}. Driver associations, delivery collectives, and labor NGOs are experimenting with tools that allow workers to log pay, trips, wait times, tips, algorithmic changes, and deactivations~\cite{hernandezEndDayAm2024,doDesigningGigWorker2024,hsiehGig2GetherDatasharingEmpower2025}. These technically modest ``folk infrastructures'' perform a crucial epistemic function: they make it possible to detect patterns of unfair treatment, evaluate the impact of algorithm updates, and substantiate claims in disputes with platforms, regulators, or policymakers~\cite{dalalLessonsWorkersInquiry2024,nagarajraoRideshareTransparencyTranslating2025,calacciAccessUnderstandingCollective2023}.

\subsubsection*{Pressures and values}
Platform labor makes enclosure structural: firms capture rich labor telemetry but disclose only thin slices, treating logs, experiments, and ranking rules as trade secrets~\cite{jarrahiAlgorithmicManagementWork2021,steinYouAreYou2023}. Openness, as infrastructuring, therefore means worker access and portability plus shared logging tools and interoperable data models that support collective analysis without relying on platform interfaces~\cite{calacciAccessUnderstandingCollective2023,hsiehGig2GetherDatasharingEmpower2025}. Exemption appears when algorithmic management is treated as outside meaningful labor oversight; platforms use legal and technical shields to deter scrutiny and normalize opacity as a competitive advantage~\cite{calacciBargainingBlackBoxDesigning2022,vincentDataLeverageFramework2021a}. Oversight emerges through independent counter-records: standardized worker logs, documented methods, and public reporting that can inform regulators, litigation, and bargaining~\cite{nagarajraoRideshareTransparencyTranslating2025,dalalLessonsWorkersInquiry2024}. Erosion is social and temporal: churn and precarity destroy institutional memory, making degradation hard to prove and leaving harms untraceable across time~\cite{hernandezEndDayAm2024,binnsNotEvenNice2025}. Ownership counters erosion through worker-governed stewardship like secure storage, continuity plans, and governance that protects contributors while enabling longitudinal evidence~\cite{calacciAccessUnderstandingCollective2023,steinYouAreYou2023}. Infrastructuring frames the observatory as a maintained capability for contestation, not a one-off dashboard.

\subsubsection*{Anticipated critiques}
Rigor is the first critique: worker-generated data may be noisy, partial, and unevenly sampled, making it easy for platforms to dismiss as ``folk statistics.'' The infrastructuring response is to treat methodological documentation as an installed base requirement: standard instruments, sampling strategies, uncertainty reporting, and triangulation with qualitative testimony and administrative records~\cite{dalalLessonsWorkersInquiry2024,calacciBargainingBlackBoxDesigning2022}. Safety is the second critique: detailed location, earnings, and interaction logs can expose contributors to retaliation, surveillance, or subpoenas. This critique demands that access tiers and data minimization be core components: consent, encryption, role-based access, and governance rules that prevent monetization and reduce compelled disclosure risk~\cite{doDesigningGigWorker2024,steinYouAreYou2023}. Leverage is the third critique: observatories may add unpaid labor without shifting power if findings do not translate into remedies. Infrastructuring answers by specifying escalation pathways up front—complaints, bargaining demands, regulatory reporting, and policy workflows—so evidence reliably triggers institutional response~\cite{vincentDataLeverageFramework2021a,calacciOrganizingEndEmployment2022,nagarajraoRideshareTransparencyTranslating2025}. The success criterion is durable, contestable evidence that can be repeatedly mobilized toward remedy, not simply visibility.

\subsubsection*{Lineages and conversations}
Several precedents supply the most directly relevant infrastructuring repertoire. Law is central because observatories gain leverage when evidence maps to enforceable rights and procedural templates for escalation pathways, especially in labor- and employment-adjacent regimes where collective data governance can be articulated as a basis for due process and remedy~\cite{calacciAccessUnderstandingCollective2023}. Accounting contributes the logic of auditability: documentation, controls, and verification routines that make counter-records credible despite platform claims of proprietary opacity~\cite{calacciBargainingBlackBoxDesigning2022,binnsNotEvenNice2025}. From adjacent conversations, worker-centered HCI and socio-technical inquiry keep power asymmetries and surveillance risks legible while developing participatory methods for building worker-facing tools and institutions under platform hegemony~\cite{steinYouAreYou2023,zhangStakeholderCenteredAIDesign2023}. Civic tech can contribute practical delivery capacity (secure tools, usability, community support), but infrastructuring insists that tool-building be anchored in stewardship and institutional pathways to remedy~\cite{hsiehGig2GetherDatasharingEmpower2025,doDesigningGigWorker2024}.

\subsubsection*{Infrastructuring synthesis}
A public interest data infrastructuring for worker observatories could commit to converting scattered experiences like pay fluctuations, opaque re-rankings, deactivations, and rule changes into a durable counter-record that can travel across time, across apps, and into arenas where remedies are possible~\cite{calacciOrganizingEndEmployment2022,steinYouAreYou2023}. Instead of imagining the observatory as a one-off dashboard, this approach builds a maintained evidentiary commons: a shared, worker-governed memory of how algorithmic management operates in practice, especially as platforms change policies, run experiments, or quietly degrade conditions~\cite{binnsNotEvenNice2025}.

Operationally, this looks like lightweight tools that fit into workers’ lives, paired with governance strong enough to protect them. Data collection should be designed to be ``safe by default'': minimal necessary detail, usable encryption, and role-based access are the conditions of participation~\cite{hernandezEndDayAm2024,doDesigningGigWorker2024}. Tiered visibility matters here: public-facing summaries can support collective organizing and media attention, while sensitive traces are held by accountable intermediaries like labor organizations, worker centers, legal aid groups, or trusted researchers who can translate evidence into enforceable claims~\cite{calacciAccessUnderstandingCollective2023,vincentDataLeverageFramework2021a}. Infrastructuring also means treating methodological credibility as part of the public interest: documentation standards and uncertainty reporting help prevent platforms from dismissing worker evidence as anecdote~\cite{dalalLessonsWorkersInquiry2024}. Most importantly, this infrastructuring is organized around leverage. The observatory is designed to feed escalation pathways—complaints, bargaining proposals, litigation strategies, and regulatory reporting—so that evidence reliably triggers institutional response and can be translated into policy-facing artifacts when needed~\cite{nagarajraoRideshareTransparencyTranslating2025,zhangDataProbesBoundary2024}. A successful worker observatory is one that outlasts churn, resists retaliation, and repeatedly converts lived experience into contestable proof and collective power, narrowing the informational asymmetries that platforms rely on to govern without accountability.


\subsection*{Case 3: Sustaining archives in the Anthropocene}

The Anthropocene is often narrated through energy transitions, climate adaptation infrastructure, and geopolitical volatility. Less visible, but equally consequential, is its challenge to how societies remember and govern. Climate disruption places extraordinary stress on the socio-technical systems that generate durable public memory: statistical agencies, environmental monitoring networks, scientific repositories, municipal records, community archives, and the data centers and power systems that keep them operational. Archives are physical infrastructures susceptible to compound shocks—heat, flood, fire, and chronic under-maintenance—at the same time that memory institutions (including libraries) are being asked to operate under widening resource constraints and governance volatility~\cite{mazurczykAmericanArchivesClimate2018,orrClimateChangeCultural2021,sesanaClimateChangeImpacts2021}. This reframing also clarifies how preservation is indelibly political: what gets measured, which vocabularies travel with datasets, and whose testimony is treated as evidence structures which injustices can later be named and which are rendered invisible~\cite{woodMobilizingRecordsReframing2014,yeoConceptsRecord12007,carbajalCriticalDigitalArchives2021a}. Climate politics makes memory adversarial, and decisions to withdraw resources from data collection, curation, or public access can quietly erase the substrate needed to validate models, reconstruct histories, or adjudicate competing claims about responsibility and harm~\cite{mazurczykAmericanArchivesClimate2018,robinsonComeHellHigh2021}. 

In the Anthropocene, archiving becomes inseparable from the conditions for future truth and reconciliation. Truth and reconciliation processes do not begin with perfect information; they begin with contested narratives and competing claims that must be stabilized through credible evidence and public procedures~\cite{robinson-sweetTruthReconciliationArchivists2018,ghaddarSpectreArchiveTruth2016,viebachEvidenceUseArchives2022}. Environmental harm characterized by colonial extraction, racialized disparate impacts, and labor exploitation often becomes legible only decades later. If the observational record has rotted away, been privatized behind restrictive licenses, or been selectively curated by institutions with conflicts of interest, future publics will have little basis to press claims for repair, restitution, or reform~\cite{wessellKeepingArchivesWater2023,ghaddarSpectreArchiveTruth2016}. The relevance of Anthropocene archives is therefore not only scientific continuity but civic possibility: whether future generations can demonstrate what happened, to whom, and under what governance decisions. A public interest approach asks how to build archives as a durable evidence commons: resilient to climate disruption, resistant to capture, and explicitly oriented to intergenerational accountability.

\subsubsection*{Pressures and values}
Climate disruption reframes enclosure, exemption, and erosion as a crisis of public memory with direct consequences for governance and future truth and reconciliation. Enclosure in environmental data is increasingly less about outright absence than about dependency: even “available” datasets can be constrained by restrictive terms, fragile interfaces, or opaque intermediaries, leaving accountability contingent on the discretion of implicated actors~\cite{carbajalCriticalDigitalArchives2021a,blankeReassemblingDigitalArchives2024}. In response, openness as infrastructuring treats access as a long-horizon commitment: durable formats, robust documentation of context and uncertainty, and licensing that preserves reuse while maintaining the linkability needed to assemble evidence across observational series, filings, disclosures, and community measurements~\cite{tanseyArchivalAdaptationClimate2015,paschalidouStrongSustainabilityFramework2022a,orrClimateChangeCultural2021}.

Exemption intensifies the stakes by enabling selective disclosure and opaque assumptions that institutionalize plausible deniability; oversight therefore requires independent custodianship capable of verification, maintaining parallel archives, and detecting disappearance or manipulation in adversarial settings~\cite{viebachEvidenceUseArchives2022,ghaddarSpectreArchiveTruth2016}. Finally, erosion compounds physical and institutional fragility as climate risks to repositories intersect with budget cuts, portal shutdowns, format obsolescence, and vendor churn. Ownership must be understood as stewardship: distributed redundancy, migration planning, interoperability, and durable governance that sustains interpretability across ``infrastructure time''~\cite{rydenArchivistsTimeConceptions2019,danielaagostinhoBigDataTime2016,wessellKeepingArchivesWater2023}. Together, infrastructuring frames Anthropocene archives as a resilient installed base for intergenerational accountability rather than a static repository of files.

\subsubsection*{Anticipated critiques}
Three critiques sharpen what Anthropocene archives must become to be credible and just. The feasibility critique notes that long-horizon stewardship is costly and complex amid underfunded, fragmented archival ecosystems; an infrastructuring response rejects a single “grand archive” in favor of federated, interoperable networks of libraries, universities, agencies, and community archives sharing tools and explicit maintenance commitments~\cite{tanseyArchivalAdaptationClimate2015,mazurczykAmericanArchivesClimate2018,robinsonComeHellHigh2021}. In this view, distributed redundancy is not inefficiency but resilience: a continuity backbone that can absorb shocks, persist through political turnover, and prevent single points of failure from becoming single points of erasure. The representation critique warns that archives can reproduce hierarchy through canon formation and dominant vocabularies that marginalize local, Indigenous, and community knowledge. Infrastructuring therefore treats provenance, categorization, and documentation as participatory governance sites, preserving context and standpoint and enabling plural knowledge systems to remain legible without being flattened into a single technocratic ontology~\cite{woodMobilizingRecordsReframing2014,carbajalCriticalDigitalArchives2021a}. The co-optation critique emphasizes that “trust” can be reputational cover and partnerships can create enclosure-by-partnership or state-centered control. Infrastructuring responds by embedding countervailing power—transparent and auditable governance, independence from dominant data owners, safeguards against restrictive licensing, and escalation pathways to contest exclusion, manipulation, or disappearance—so the archive is accountable not only for what it stores but for how it governs~\cite{blankeReassemblingDigitalArchives2024,viebachEvidenceUseArchives2022}.

\subsubsection*{Lineages and conversations}
Anthropocene archiving in anticipation of future truth and reconciliation efforts draws on complementary public-interest lineages that together specify how climate memory can become durable, credible, and actionable. Engineering contributes high-reliability stewardship sensibilities—redundancy, monitoring, incident learning, and maintenance regimes—that treat breakdown as inevitable and continuity as something designed, resourced, and governed under stress~\cite{wessellKeepingArchivesWater2023,mazurczykAmericanArchivesClimate2018}. Accounting and auditing add verification logics essential for climate claims-making, emphasizing documentation and routinized scrutiny so records remain credible and assumptions and uncertainty can be evaluated rather than asserted~\cite{paschalidouStrongSustainabilityFramework2022a,orrClimateChangeCultural2021}. Law and transitional-justice practice provide enforceable duties and procedural pathways that connect evidence to remedy—retention obligations, access rights, discovery rules, and due-process protections—while also structuring tiered access and protected disclosure to manage privacy and sovereignty concerns~\cite{robinson-sweetTruthReconciliationArchivists2018,ghaddarSpectreArchiveTruth2016}. Among adjacent conversations, critical archival and digital-archives traditions sharpen the political economy of memory by foregrounding how infrastructures of access, description, and platform mediation can reproduce inequality or facilitate capture~\cite{carbajalCriticalDigitalArchives2021a,woodMobilizingRecordsReframing2014}. Finally, practical climate-archival action within memory institutions (including libraries) offers near-term footholds for turning “stewardship” into budgets, workflows, and procurement decisions under climate stress~\cite{ajaniLibrariesClimateChange2025,robinsonComeHellHigh2021}.

\subsubsection*{Infrastructuring synthesis}
A public interest approach to Anthropocene archives treats climate memory as resilience infrastructure and a prerequisite for future truth and reconciliation. It shifts archiving from a back-office technical task to a civic institution with explicit obligations to future publics who will need durable evidence to reconstruct what happened, assign responsibility, and press claims for repair under conditions of physical disruption and political contestation. This reframing is also temporal: it asks what it means to design and govern evidence for deep time, when the relevant publics are not only dispersed across jurisdictions but also across generations~\cite{holtorfCulturalHeritageFuture2020,rahm-skagebyHCIDeepTime2022,hanuschDeeptimeOrganizationsLearning2020}. The core strategy is a federated “continuity backbone” spanning libraries, universities, scientific repositories, agencies, community archives, and watchdog organizations, where distributed mirrors reduce vulnerability to disaster and austerity while interoperable standards prevent fragmentation from becoming unintelligibility. 

Infrastructuring makes openness durable by insisting on practical linkability and interpretability: stable identifiers that allow records to be joined across systems, and provenance practices that preserve measurement context, assumptions, and uncertainty so that future audits can distinguish disagreement from deception~\cite{yeoConceptsRecord12007,josiasUnderstandingArchivesFeature2011}. Within this backbone, oversight and ownership are treated as governance capacities rather than optional add-ons. Oversight means independent custodianship that can operate under adversarial incentives—detecting disappearance, flagging degradation, preserving parallel copies of at-risk datasets, and verifying contested claims—while making the archive’s own rules for inclusion, exclusion, and access transparent and contestable~\cite{blankeReassemblingDigitalArchives2024,ghaddarSpectreArchiveTruth2016}. Ownership becomes stewardship with teeth: durable, collectively governed responsibilities that outlast electoral and grant cycles, coupled to access tiers and accountable intermediaries that protect communities from extraction and harm while still enabling safe scrutiny. Finally, reconciliation-ready archives require escalation pathways—procedures for challenging omissions, triggering review, and contesting category choices—so that the archive remains not only preserved but governable. The promise is not preservation for its own sake, but a maintained evidentiary commons that can support verification of climate targets, substantiation of environmental harms, and reconciliation processes that have not yet been invented~\cite{viebachEvidenceUseArchives2022,robinson-sweetTruthReconciliationArchivists2018}.


\section*{Discussion}
Public interest data infrastructuring advances a unifying framework for a rapidly changing political economy of data. Rather than treating shrinking access as an inconvenience for researchers, the paper argues that enclosure, exemption, and erosion reorganize the institutional conditions under which public knowledge is produced and used with implications for democratic oversight. In response, it proposes openness, oversight, and ownership not as aspirational values but as a theory of change: each names a civic capacity that must be rebuilt as data-driven systems increasingly function as critical infrastructure.

The paper's central theoretical contribution is to clarify infrastructuring as the mechanism that translates these values into durable accountability capacities. Infrastructuring shifts emphasis from isolated disclosures, dashboards, or episodic audits toward maintaining an installed base of records, standards, interfaces, and governance routines that makes systems observable, auditable, and contestable over time. Operationally, the installed base checklist (identifiers, provenance, retention, access tiers, escalation pathways) supplies a practical vocabulary for diagnosing failure modes and specifying design constraints. It also creates a comparative lens for evaluating adjacent efforts: initiatives that deliver short-term analytical wins can still fail the infrastructuring test if they do not institutionalize durable reference, interpretability, preservation, safe access, and routinized remedy.

A second contribution synthesizing across public interest lineages (journalism, law, accounting, planning, engineering) and adjacent conversations (public interest technology, data for good, GovTech, civic tech, critical data studies, data justice) through the installed base lens. Journalism and planning help constitute publics around evidence and keep categories contestable; law and accounting translate values into enforceable disclosure, verification, and remedy; engineering contributes high-reliability practices for governing complex systems; and critical data studies and data justice discipline the agenda by insisting that visibility regimes and documentation practices distribute burdens and benefits. These traditions and conversations highlight why coalition is structurally necessary. Contestation traditions sharpen what the installed base must be for (whose interests, which harms, what refusal looks like); accountability traditions contribute enforceable recordkeeping and escalation machinery; and delivery traditions implement interoperable systems that can persist across procurement cycles and institutional churn. 

The three case studies extend these arguments into applied implications. Mining municipal archives frames ``open government'' as subject to familiar forces of soft enclosure via unusable formats and weak retrieval, exemption via compliance-oriented disclosure without auditability, and erosion via link rot and custodial churn. Worker observatories illustrate infrastructuring under adversarial conditions, where the object of accountability is an employer: the core deliverable becomes an independent counter-record with governance strong enough to withstand retaliation and to translate evidence into remedies. Anthropocene archives push the framework into intergenerational time, treating climate-relevant memory as a resilience and justice problem whose feasibility depends on federated stewardship, participatory provenance, and anti-capture governance. Across cases, the recurring design pattern is tiered visibility with accountable intermediaries: a way to preserve scrutiny without converting openness into a new vector of extraction or surveillance.

Limitations and alternatives sharpen the agenda. First, infrastructuring is capacity intensive: local jurisdictions, labor groups, and civic institutions lack sustained funding, staff, or legal support. The paper's proposals can read as aspirational where public infrastructure is already precarious and brittle. Second, the checklist can be misapplied as a technocratic recipe: identifiers and retention can increase exposure if access tiers and governance are weak, and escalation pathways can devolve into performative ``accountability theater'' without enforceable hooks. Third, the framework does not claim that infrastructuring is the only route to reform. Open science, responsible technology toolkits, and project-based ``data for good'' efforts can still matter, especially as rapid response capacity. The paper's claim is narrower: without installed base commitments, these efforts are unlikely to survive enclosure and erosion, or to consistently counter exemption.

Future work across research, practice, and institutions. Empirically, the field needs comparative studies of ``installed base failures'' (\textit{e.g.}, where provenance collapses, where retention is strategically undermined, where escalation is blocked) and evaluations of which governance arrangements most reliably prevent capture. Methodologically, it needs robust patterns for audit-ready extraction from messy corpora, privacy-preserving observatories, and verifiable change-logging for continuously updated models. Institutionally, it needs proposals for statutory access rights, mandated research portals, retention duties, and independent audit capacities that are resourced and durable; professional pathways that treat stewardship and oversight as first-class contributions; and partnerships with libraries and archives as long-horizon custodians of public memory. The installed base checklist is intended as a shared scaffold for this program: a way to turn broad commitments into testable, maintainable accountability infrastructure.

\section*{Conclusion}
Public interest data infrastructuring is a necessary pivot in how we conceptualize data-driven accountability. Under enclosure, exemption, and erosion, democratic oversight cannot be sustained by goodwill, episodic access, or one-off interventions. Accountability is a maintained capability built and rebuilt through infrastructuring an installed base of identifiers, provenance, retention, access tiers, and escalation pathways that keep powerful systems observable, auditable, and contestable over time.  By triangulating public interest professional lineages with adjacent conversations in civic and critical scholarship, the paper locates complementary capacities and clarifies why coalition is not optional. The case studies show how these repertoires can be mobilized in concrete settings: municipal archives as a test of usable openness, worker observatories as counter-records under adversarial control, and Anthropocene archives as intergenerational stewardship under climate disruption. If societies want data and models to function as part of the democratic commons rather than as instruments of private enclosure and regulatory exemption, they must invest in durable accountability infrastructure as a set of technical, legal, and institutional commitments.

\begin{acks}
Thanks for all the fish.
\end{acks}

\bibliographystyle{ACM-Reference-Format}
\bibliography{bibliography}

\end{document}